\documentclass[11pt]{article}

\usepackage[margin=1in]{geometry}
\usepackage[T1]{fontenc}
\usepackage[utf8]{inputenc}
\usepackage{lmodern}
\usepackage{microtype}
\usepackage{amsmath,amssymb,amsthm,mathtools}
\usepackage{enumitem}
\usepackage{booktabs}
\usepackage{algorithm}
\usepackage{algpseudocode}
\usepackage{xurl}
\usepackage[hidelinks]{hyperref}
\setlength\emergencystretch{2em}

\newtheorem{theorem}{Theorem}
\newtheorem{lemma}{Lemma}
\newtheorem{definition}{Definition}
\newtheorem{remark}{Remark}

\newcommand{\TV}{\mathrm{TV}}

\title{Proof-Carrying Budgets for Anonymous DHTs\\
\large Dual certificates and tiny verifiers for padding, delegation, and congestion contracts}
\author{}
\date{Draft (2026-03-22; archive v502)}

\begin{document}
\maketitle

\begin{abstract}
Anonymous DHT designs increasingly resemble \emph{compiled systems}: a lookup compiler chooses padding/deadlines, delegation lengths, retry rules, and congestion knobs, then an ``accountant'' reports a privacy budget.
If those budgets are meaningful, they should be reproducible by parties who do \emph{not} trust the compiler.

We present a proof-carrying pattern for privacy budgets.
A deployment publishes a compact \emph{verifier bundle}: (i) a concrete plan (primal witness), (ii) a dual witness that upper-bounds the privacy objective, and (iii) deterministic evaluator rules that let a tiny checker replay the declared bound.
The checker verifies feasibility, evaluates the primal/dual arithmetic, and confirms the certificate inequality without rerunning large experiments.

This paper therefore focuses on the narrow \emph{verifier-contract interface}: what replayable object a budget-carrying anonymous-DHT paper can hand upward once the math is fixed.
It is intentionally the optional verifier-contract companion beneath Boss Fight~A / Boss Fight~B rather than a second accountant, evaluator, or release paper.
ABOM/OINL claim surfaces, notarized attempt logs, and signed release receipts are treated as imported packaging layers owned by adjacent papers in this archive.
\end{abstract}

\paragraph{Scope.}
We isolate and formalize \emph{Proof-Carrying Budgets for Anonymous DHTs} as the replayable verifier-contract companion in the Boss Fight stack: once Boss Fight~A has fixed the repeated-use accountant and Boss Fight~B has fixed the anonymous-DHT knob surface, this note exports the minimal object a checker can replay without trusting the compiler.

\paragraph{Imports.}
We treat (i) padding duals/transport certificates as owned by the published Mathematics-era dual-certificate paper on optimal poset-feasible padding and transport,
(ii) stop\,--\,time normal forms as owned by the published Mathematics-era stop-time addendum,
(iii) repeated-use horizon accounting and knob selection as owned by Boss Fight~A / Boss Fight~B,
(iv) optional table-heavy operator evidence as owned by the Boss Fight~B addendum, and
(v) notarization / claim packaging / release lineage as owned by Eval~1, ABOM/OINL, and Release~A.\cite{series:bossfighta,series:bossfightb,series:bossfightBaddendum,series:evalnotary,series:abomoinl,series:releaseproperty}

\paragraph{Exports.}
A proof-carrying \emph{verifier bundle}: declared plan schema, declared dual schema, deterministic primal/dual evaluators, the certificate inequality $\mathrm{PrimalValue}(\pi)\le \mathrm{DualValue}(\lambda)$, and a short checker that reproduces the published bound from $(\pi,\lambda)$ alone.
This paper does \emph{not} own operator-facing knob sweeps, notarized attempt logs, ABOM/OINL claim diffs, or signed release receipts; it exports only the narrow replay object those later layers commit to.\cite{series:bossfightBaddendum,series:evalnotary,series:abomoinl,series:releaseproperty}

\paragraph{Stable handoff.}
The clean handoff to adjacent papers is:
\emph{declared knob surface / plan choice $\rightarrow$ optional evidence tables $\rightarrow$ verifier bundle $(\pi,\lambda)$ $\rightarrow$ notarized evaluation bundle $\rightarrow$ ABOM/OINL claim surface $\rightarrow$ signed release receipt}.
Boss Fight~B chooses the anonymous-DHT knobs, the addendum may publish compact operator tables for those knobs, this paper turns the declared choices into a replayable verifier object, Eval~1 notarizes attempted evaluations of that object, and ABOM/OINL plus Release~A bind the resulting claim to a public release.\cite{series:bossfightb,series:bossfightBaddendum,series:evalnotary,series:abomoinl,series:releaseproperty}
Relative to the Boss Fight stack, Paper~A is the repeated-use theorem/accountant entrypoint, Paper~B is the practical dial sheet, the addendum is optional supporting evidence, and this note is only the verifier-contract companion.

\paragraph{Later-terse-paper citation posture.}
When a later short paper only needs to answer \emph{which replayable verifier contract was declared for this anonymous-DHT budget?}, it should stop at the minimal public verifier-bundle card below plus the tiny checker drill.
It should widen only when the live issue is knob selection (Boss Fight~B), optional operator evidence (the addendum), evaluator notarization (Eval~1), or release-lineage packaging (ABOM/OINL and Release~A).\cite{series:bossfightb,series:bossfightBaddendum,series:evalnotary,series:abomoinl,series:releaseproperty}

\paragraph{Artifact policy.}
This archive is source-only; supporting artifacts are available on request. See \cite{series:artifactpolicy}.

\section{The proof-carrying template}
Many ``privacy budgets'' in this series can be written as optimization problems:
choose a mechanism $\pi$ (a padding schedule, a randomized delegation rule, a congestion calendar) to minimize some cost subject to a privacy constraint, or dually maximize privacy subject to a cost envelope.
When the budget is stated in \emph{TV form} or another convex divergence, these problems admit duals whose feasibility can be checked locally.

\begin{definition}[Certificate (informal)]
A \emph{budget certificate} is a tuple $(\pi,\lambda)$ where $\pi$ is a concrete mechanism instance (a plan) and $\lambda$ is a dual witness such that:
(i) $\pi$ satisfies the stated feasibility constraints;
(ii) $\lambda$ is dual-feasible; and
(iii) the dual objective evaluated at $\lambda$ upper-bounds the privacy objective of $\pi$.
\end{definition}

\begin{lemma}[What the verifier checks]
If the verifier can (a) re-evaluate the primal objective/constraints and (b) re-evaluate the dual objective/constraints from $(\pi,\lambda)$, then the verifier can reproduce the published privacy bound with no access to private traces.
\end{lemma}

\section{A minimal verifier interface}
The verifier interface is intentionally small; it should fit in a few hundred lines.
\begin{algorithm}[h]
\caption{Toy verifier interface for a budget certificate}
\begin{algorithmic}[1]
\Procedure{Verify}{$\mathcal P,\mathcal D,\pi,\lambda,B$}
  \State \textbf{check} $\pi\in\mathcal P$ (schema + feasibility)
  \State \textbf{check} $\lambda\in\mathcal D$ (schema + dual feasibility)
  \State $p\gets \mathrm{PrimalValue}(\pi)$
  \State $d\gets \mathrm{DualValue}(\lambda)$
  \State \textbf{check} $p\le d$ \Comment{weak duality / certificate inequality}
  \State \textbf{check} $d\le B$ \Comment{published bound}
  \State \Return \textsf{ok}
\EndProcedure
\end{algorithmic}
\end{algorithm}

\paragraph{What counts as ``primal'' here?}
For padding, $\pi$ can be a discrete distribution over stop times (or a mixture over deadline normal forms).
For committee contact privacy, $\pi$ can be a randomized contact schedule (with public randomness commitments).
For congestion, $\pi$ can be the compiler knob manifest plus the accountant transcript.

\section{Exported object and earliest sufficient stop}
The public object owned by this paper is intentionally narrow:
\begin{center}
\fbox{\parbox{0.92\linewidth}{\centering
\textbf{verifier bundle} $=$ declared plan schema $+$ declared dual schema $+$ replayable evaluators $+$ certificate inequality check}}
\end{center}
An auditor who can parse that bundle, recompute the primal/dual arithmetic, and confirm the declared bound has already extracted everything this paper promises.
That is the earliest sufficient stop.
Questions about whether the bundle was accompanied by knob-sweep evidence, retried, selectively disclosed, or attached to the correct release are explicitly deferred to the Boss Fight~B addendum, Eval~1, and ABOM/OINL + Release~A.\cite{series:bossfightBaddendum,series:evalnotary,series:abomoinl,series:releaseproperty}

\paragraph{A minimal public verifier-bundle card.}
Table~\ref{tab:minimal-verifier-card} is the smallest public packet later terse papers should need when the live question is simply \emph{which replayable checker contract was declared for this budget-carrying anonymous-DHT plan?}

\begin{table}[t]
\centering
\small
\begin{tabular}{@{}p{0.31\linewidth}p{0.60\linewidth}@{}}
\toprule
\textbf{Public row} & \textbf{Pinned value}\\
\midrule
Declared knob surface & Boss Fight~B knob card \texttt{bf-knobs-v1}: delegation length $4$, dummy rate $0.15$, retry cap $1$, congestion pad factor $0.20$ \\
Plan schema / plan ID & \texttt{bf-plan-v1} / \texttt{plan-7} \\
Dual schema / witness ID & \texttt{bf-dual-v1} / \texttt{dual-7} \\
Evaluator version & deterministic checker \texttt{bf-checker-v1} for primal and dual replay \\
Replayed values & $p=\mathrm{PrimalValue}(\pi)=0.041$ bits, $d=\mathrm{DualValue}(\lambda)=0.047$ bits \\
Published cap & $B=0.050$ bits for the declared repeated-use window \\
Verifier verdict & accept iff $p\le d\le B$; here $0.041\le 0.047\le 0.050$ \\
Downstream owner & Eval~1 / ABOM--OINL / Release~A for notarization, claim packaging, and lineage \\
\bottomrule
\end{tabular}
\caption{A minimal public verifier-bundle card. This is the non-decorative public answer later terse papers can quote directly when the live issue is the declared replayable checker contract rather than knob selection, attempt notarization, or release packaging.}
\label{tab:minimal-verifier-card}
\end{table}

This card is intentionally narrow: it pins the declared plan/dual schemas, the deterministic checker version, and one concrete $p\le d\le B$ verdict, but nothing else.
Boss Fight~B still owns how the knob surface was chosen, the addendum still owns any optional supporting tables, and Eval~1 plus ABOM/OINL plus Release~A still own attempted-evaluation disclosure and public lineage.\cite{series:bossfightb,series:bossfightBaddendum,series:evalnotary,series:abomoinl,series:releaseproperty}

\paragraph{One tiny checker drill.}
With the card values above, the checker recomputes
\[
p=0.041,\qquad d=0.047,\qquad B=0.050,
\]
and accepts because
\[
0.041\le 0.047\le 0.050.
\]
A later terse paper can quote the card and this inequality without reopening the knob compiler, evaluator-attempt ledger, or release stack.

\paragraph{A fail-closed verifier-bundle cutover.}
The smallest public cutover after a verifier bundle exists is not merely numerical.
Suppose a later deployment reuses the same Boss Fight horizon card and the same anonymous-DHT knob card, and its replay still lands at
\[
p=0.041,\qquad d=0.047,\qquad B=0.050.
\]
If it still declares \texttt{bf-plan-v1}, \texttt{bf-dual-v1}, and \texttt{bf-checker-v1}, then the old verifier-bundle card remains the right public reference.
But if the deployment switches to \texttt{bf-checker-v2} (or changes the declared plan/dual schema IDs) while keeping the same numbers, the old card no longer names the actual replay contract: the bundle must be re-issued under the new checker/schema identifiers even before any higher-layer notarization or release packaging is discussed.
By contrast, if the live change is the published cap itself---for example because the repeated-use horizon or the declared knob surface changed and the bundle now needs $B=0.060$ rather than $0.050$---then the issue has already widened out of this paper and back into Boss Fight~A or Boss Fight~B rather than being a mere verifier-bundle refresh.\cite{series:bossfighta,series:bossfightb}

\paragraph{Minimal reopen ladder.}
Later terse papers should therefore stop at the verifier bundle only while the live question is \emph{which declared plan/dual/checker contract produced this replayable $p\le d\le B$ verdict?}
The first widening owner is then explicit:
\begin{itemize}[leftmargin=*]
\item reopen \textbf{Boss Fight~A} if the repeated-use horizon / published cap itself changed,\cite{series:bossfighta}
\item reopen \textbf{Boss Fight~B} if the anonymous-DHT knob surface or plan family changed,\cite{series:bossfightb}
\item reopen the \textbf{Boss Fight~B addendum} if the question is now about optional evidence tables or cautionary operator slices rather than the replayable checker contract,\cite{series:bossfightBaddendum}
\item reopen \textbf{Eval~1} if the live issue is attempted-evaluation notarization, retry discipline, public randomness, or the gap-free attempt prefix,\cite{series:evalnotary}
\item reopen \textbf{ABOM/OINL} if the live issue is the signed claim-object surface or sameness of the public anonymity claim,\cite{series:abomoinl}
\item and reopen \textbf{Release~A} if the live issue is lineage-tagged publication, signed receipts, or fail-closed release gating.\cite{series:releaseproperty}
\end{itemize}
This keeps this note narrow: it owns the replayable verifier contract and its checker-level cutover, but not the upstream accountant / knob choices and not the downstream notarization / claim-packaging / release layers.

\section{Handoffs beyond this paper}
Proof objects must still be publishable, but the outer layers are owned elsewhere.
We therefore recommend the following layering (owned by other papers):
\begin{itemize}[leftmargin=*]
\item \textbf{Boss Fight~B / domain papers:} choose the declared knob surface and plan family to which the verifier bundle applies.\cite{series:bossfightb}
\item \textbf{Boss Fight~B addendum / evidence companions:} optionally publish compact operator tables, pattern maps, and cautionary slices for the declared knob family without changing the verifier contract itself.\cite{series:bossfightBaddendum}
\item \textbf{Eval~1:} notarize attempted evaluations of the verifier bundle, including retry policy, public randomness binding, and the gap-free attempt prefix.\cite{series:evalnotary}
\item \textbf{ABOM/OINL + Release~A:} commit the verifier bundle by digest, summarize the resulting claim surface, and bind it to a release lineage / signed receipt.\cite{series:abomoinl,series:releaseproperty}
\item \textbf{Specialized domain certificates:} congestion budgets add log inclusion/consistency proofs and monitoring receipts (W-Congestion-EQ) on top of the same verifier skeleton.\cite{series:wcongeq}
\end{itemize}

\section{Limitations}
A dual certificate validates arithmetic and consistency, not ground truth.
If the underlying model misses a side channel, a perfect certificate can still describe the wrong world.
We therefore treat proof-carrying budgets as one layer in a larger publication pipeline that also includes state-dependent threat windows, closed-view auditing, retry-notarization, and release-lineage checks in later papers.

\section{Takeaways}
\begin{itemize}[leftmargin=*]
\item \textbf{Proof-carrying budget pattern:} publish a concrete plan $\pi$ and a dual witness $\lambda$ so third parties can re-check the stated TV/privacy bound by weak duality.
\item \textbf{Own only the verifier bundle:} this paper stops once an auditor can parse the declared schemas, replay the primal/dual arithmetic, and confirm the bound from $(\pi,\lambda)$ alone.
\item \textbf{Keep outer layers separate:} Boss Fight~B chooses knobs, the addendum optionally ships extra operator tables, Eval~1 notarizes attempted evaluations, ABOM/OINL summarizes the claim surface, and Release~A binds that surface to a release lineage.
\item \textbf{Limits:} certificates validate arithmetic and consistency, not whether the model omitted a side channel---pair with threat-window analysis, auditing, and retry discipline.
\end{itemize}

\begin{thebibliography}{99}\small

\bibitem{series:artifactpolicy}
Anonymous.
\newblock Artifact and Disclosure Policy for the One-True-Archive (Source-Only Public Release).
\newblock Series synthesis note (Synthesis~6), 2026. (This archive.)

\bibitem{series:bossfighta}
Anonymous.
\newblock Boss Fight Budgets: Maximal Leakage as a Repeated-Use Anonymity Contract.
\newblock Boss Fight series Paper~A, 2026. (This archive.)

\bibitem{series:bossfightb}
Anonymous.
\newblock AnonDHT in the Boss Fight Regime: Compiling a MaxL Budget into Lookup Knobs.
\newblock Boss Fight series Paper~B, 2026. (This archive.)

\bibitem{series:bossfightBaddendum}
Anonymous.
\newblock AnonDHT in the Boss Fight Regime (Addendum): Evidence Tables and Pattern Map.
\newblock Boss Fight series Paper~B addendum, 2026. (This archive.)

\bibitem{series:evalnotary}
Anonymous.
\newblock Notarized Anonymity Certificates Under Retries: Grinding-Resistant Evaluation via Transparency Logs, Public Randomness, and $\delta$-Spending.
\newblock Evaluation series Paper~1, 2026. (This archive.)

\bibitem{series:abomoinl}
Anonymous.
\newblock ABOM and OINL for Anonymity Claims: Proof-Carrying Odds-Inflation Budgets Without Publishing Full Artifacts.
\newblock Anonymity series Paper~C, 2026. (This archive.)

\bibitem{series:releaseproperty}
Anonymous.
\newblock Privacy as a Release Property: Proof-Carrying Anonymity Receipts for Anonymous DHTs.
\newblock Release and destination Paper~A, 2026. (This archive.)

\bibitem{series:wcongeq}
Anonymous.
\newblock W-Congestion-EQ: Verifiable Congestion Budgets --- Congestion-Specific Verifier Bundles for Anonymous Overlays.
\newblock Congestion series Paper~3, 2026. (This archive.)

\end{thebibliography}

\end{document}
